\section{Pancreatic model controls: clinical covariates do not replace mutation-context prediction}
The supporting pancreatic study predicts nonsynonymous CDKN2A mutation status from the rest of the tumor's mutation landscape. It is not a model of treatment benefit, diagnosis or prospective prognosis. Earlier GNN attribution rankings failed to replicate across splits, and the sixteen-split ensemble did not make them stable. Those negatives remain unchanged. A new frozen local plan instead asks a smaller model-validation question: how much held-out discrimination can simple clinical versus target-excluded mutation-context features supply under a fixed logistic estimator?

Inspecting the implementation also revealed a label error. The clinical feature builder returns age and an is-male indicator. The covariate-permutation script called its second feature ``smoker'', although it never loaded a smoking variable. We correct that script label to \texttt{is\_male}. No numeric output generated from that feature is relabeled as smoking evidence. Missing sex is still encoded zero by the historical source convention; zero therefore is not a verified female label for every patient. The new audit retains this limitation and does not infer demographics from unmeasured variables.

The local TCGA-PAAD files contain 184 samples and 184 distinct patient barcode prefixes, with 35 CDKN2A-positive samples. This equality means patient grouping does not remove duplicate samples in this particular pool, but grouping remains explicit rather than assumed. Stratified shuffled patient splits use fixed seeds 17, 29 and 43. Each split has 138 training and 46 test patients, nine positive in each test set. There are no missing ages in this pool, so imputation does not change any record here. The age median is nevertheless learned only from the training set (65.5, 65 and 65 years respectively) and used to impute missing test ages. Scaling is also fitted on training data only. The earlier feature builder imputed age from the entire pool before splitting; the new workflow avoids that transductive preprocessing without silently rewriting the old experiment.

Three preselected views use the same held-out patients: age plus is-male; target-excluded mutation burden, truncating fraction, number of mutated genes and six pathway indicators; and their concatenation. CDKN2A itself is excluded from mutation counts and pathway inputs. A balanced, L2-penalized logistic regression with $C=1$ and at most 2,000 iterations is fixed for all views and seeds. There is no hyperparameter search or selection of the best split. This is not a retrained GNN, so its scores cannot be compared with a GNN trained on different splits as a model-superiority test.

\begin{table}[ht]\centering\small
\caption{Held-out CDKN2A status AUC for three fixed patient splits, nine positives among 46 test patients in each. All splits are retained. They overlap and are not independent clinical replications.}
\begin{tabular}{lrrr}\hline
View & seed 17 & seed 29 & seed 43\\\hline
Age and is-male & .5255 & .5465 & .4550\\
Burden and pathways & .6967 & .7868 & .7868\\
Combined & .7087 & .7748 & .7628\\\hline
\end{tabular}\end{table}

Clinical age/sex alone is near chance and below chance in one tested split. Mutation-context controls achieve higher AUC in all three, but adding the clinical variables raises the seed-17 AUC and lowers the other two. The result does not support a universal claim that more clinical inputs improve this predictor. It also cannot identify why mutation context predicts status: co-mutation structure, burden, cohort composition and ascertainment may all contribute. Correlated pathway and burden indicators do not establish causal biological pathways or justify selecting a therapy.

Average precision and probability-quality measures are recorded separately rather than treating discrimination as calibration. Mutation-context average precision is .3150, .4403 and .4305; combined-view values are .3251, .3996 and .4501. The target prevalence in each test fold is $9/46=.19565$, so these must be read against class imbalance. Combined Brier scores are .2440, .1959 and .1836, and log losses are .6678, .5397 and .5226. Class-balanced logistic fitting changes the training loss weights and does not guarantee prevalence-calibrated probabilities. No decision threshold, clinical utility curve or calibrated deployment model is supplied.

With only nine positive patients per test fold, score changes can arise from a small number of reordered pairs. The three splits reuse the same cohort and cannot be averaged into an independent-patient confidence interval. We provide no significance claim or external-validity estimate. Mutation-context inputs are retrospective features from the same assay family as the target. Excluding the target mutation avoids direct identity leakage but does not turn this task into an independent test of biological mechanism. Broader external-cohort validation and prospective outcome data remain missing.

\path{scripts/pancreatic_control_plan.json} records the frozen estimator, views and split seeds; \path{scripts/pancreatic_control_audit.py} executes the patient-grouped training-only preprocessing. \path{results/pancreatic_control_audit.json} retains patient partitions, missing-age count, mutation-source hash, every AUC/AP/Brier/log-loss record and the corrected clinical-feature meaning. The sample/patient prediction ledger is \path{results/pancreatic_control_predictions.jsonl}. Tests check patient disjointness, exact test counts, target exclusion, ledger hashes and metric recalculation. The historical GNN and attribution files are untouched. This study adds actual control-model results and a source-label correction, not a new pancreatic driver, biomarker or treatment recommendation.
